\documentclass{article}
\usepackage{lmodern}
\usepackage{amsmath}
\usepackage{listings}
\usepackage{fullpage}
\usepackage{parskip}
\usepackage{xcolor}
\lstset{
	basicstyle=\ttfamily,
	columns=fullflexible,
	frame=single,
	breaklines=true,
	postbreak=\mbox{\textcolor{red}{$\hookrightarrow$}\space},
}
\begin{document}
\title{Experiment `p\_6\_1\_a\_process\_array\_even\_indexed\_array' Results}
\author{\today}
\date{}
\maketitle
\textbf{Experiment outcome: }SUCCESS
\\\textbf{Bad responses: }0
\\\textbf{Responses containing}~\texttt{assume}~\textbf{: }0
\\\textbf{Responses containing}~\texttt{expect}~\textbf{: }0
\\\textbf{Resolution attempts: }1
\\\textbf{Hard fails (resolution): }0
\\\textbf{Soft fails (resolution): }0
\\\textbf{Verification attempts: }1
\section*{Problem Specification}
\textbf{Problem name: }p\_6\_1\_a\_process\_array\_even\_indexed\_array
\\\textbf{Natural language statement: }null
\\\textbf{Method signature: }p\_6\_1\_a\_process\_array\_even\_indexed\_array(inputs: array<int>) returns (even\_indexed: array<int>)
\subsection*{Ensures}
\begin{itemize}
\item \texttt{even\_indexed[..] == [inputs[0], inputs[2], inputs[4], inputs[6], inputs[8]]}
\end{itemize}
\subsection*{Requires}
\begin{itemize}
\item \texttt{inputs.Length == 10}
\end{itemize}
\clearpage
\section*{GenAI interactions}
Below you will find all interactions between the `user' (program) and the `assistant' (OpenAI).
\subsection*{Program $\rightarrow$ GenAI}
\begin{lstlisting}
You are given the following task to perform in Dafny:

Write a method to satisfy the below requirements.

The signature should be:

p_6_1_a_process_array_even_indexed_array(inputs: array<int>) returns (even_indexed: array<int>)

The method should respect the following contract:

requires inputs.Length == 10, ensures even_indexed[..] == [inputs[0], inputs[2], inputs[4], inputs[6], inputs[8]]

Produce and show only the Dafny body of this method, including the curly braces that surround it. Do not show the signature nor contract. You must not use 'assume' nor 'expect' anywhere in your code.
\end{lstlisting}
\subsection*{GenAI $\rightarrow$ Program}
\textbf{System fingerprint: }null
\\\textbf{ID: }chatcmpl-DnNSnqesNUFb2toDKHIIQ5t1lbnWw
\\\textbf{Model: }gpt-5.5-2026-04-23
\\\textbf{Created at: }1780660637
\\\textbf{Finish reason: }stop
\\\textbf{Usage: }318 tokens in, and 614 tokens out
\begin{lstlisting}
{
  even_indexed := new int[5];
  even_indexed[0] := inputs[0];
  even_indexed[1] := inputs[2];
  even_indexed[2] := inputs[4];
  even_indexed[3] := inputs[6];
  even_indexed[4] := inputs[8];
}
\end{lstlisting}
\section*{Final Program}
\begin{lstlisting}
method p_6_1_a_process_array_even_indexed_array(inputs: array<int>) returns (even_indexed: array<int>)
	requires inputs.Length == 10
	ensures even_indexed[..] == [inputs[0], inputs[2], inputs[4], inputs[6], inputs[8]]
{
  even_indexed := new int[5];
  even_indexed[0] := inputs[0];
  even_indexed[1] := inputs[2];
  even_indexed[2] := inputs[4];
  even_indexed[3] := inputs[6];
  even_indexed[4] := inputs[8];
}
\end{lstlisting}
\section*{Total Token Usage}
\textbf{Input tokens: }318
\\\textbf{Output tokens: }614
\\\textbf{Reasoning tokens: }512
\\\textbf{Sum of `total tokens': }932
\section*{Experiment Timings}
\textbf{Overall Experiment} started at 1780660636153, ended at 1780660688331, lasting 52178ms (52.18 seconds)
\\\textbf{Iteration \#1} started at 1780660636154, ended at 1780660688331, lasting 52177ms (52.18 seconds)
\end{document}
